\section{Introduction}
\label{sec:introduction}

Let $\Sym^n$ denote the real symmetric $n\times n$ matrices. A
\emph{spectrahedron} is a set defined by a finite affine linear matrix
inequality; a \emph{spectrahedral shadow} is a linear projection of such
a set. Thus auxiliary variables may make a semidefinite representation
possible even when no representation exists in the original variables.
We study this distinction for one explicit hyperbolicity cone.

For a homogeneous real polynomial $p$ with $p(e)>0$, hyperbolicity with
respect to $e$ means that $t\mapsto p(te-x)$ has only real roots for
every $x$. We use the closed-cone convention
\[
 K(p,e)=\{x:\text{every root of }t\mapsto p(te-x)
                    \text{ is nonnegative}\}.
\]
In particular, zero roots are allowed.

The convexity of hyperbolicity cones goes back to
G\aa rding~\cite{garding1959}. In three variables, the
Helton--Vinnikov theorem yields definite symmetric determinantal
representations, and hence spectrahedral
cones~\cite{helton-vinnikov2007,lewis-parrilo-ramana2005}.
Allowing auxiliary variables leads to a broader representation question.
Netzer and Sanyal proved semidefinite liftability when every nonzero
boundary point is a smooth point of the defining hyperbolic
polynomial~\cite{netzer-sanyal2015}. More recently, Scheiderer obtained
second-order cone lifts under a Nash-smooth boundary
hypothesis~\cite{scheiderer2025}. Our result instead concerns an explicit
cone and gives its lift by a bounded-degree algebraic construction.

\subsection{The cone and the main result}
For $y=(y_1,y_2,y_3)\in\R^3$, define
\begin{equation}\label{eq:Q}
 Q(y)=\begin{pmatrix}
 y_1^2+y_2^2&-y_1y_2&-y_1y_3\\
 -y_1y_2&y_2^2+y_3^2&-y_2y_3\\
 -y_1y_3&-y_2y_3&y_3^2+y_1^2
 \end{pmatrix}.
\end{equation}
The associated biquadratic form $z^TQ(y)z$ is the coefficient-one
Choi--Lam form~\cite[Section~4]{choi-lam1977}.
For $a\in\R$, $r\in\R^3$, and $T\in\Sym^3$, set
\begin{equation}\label{eq:Phi}
 \Phi_y\!\begin{pmatrix}a&r^T\\r&T\end{pmatrix}
 =\begin{pmatrix}
 \tr(Q(y)T)&-r^TQ(y)\\
 -Q(y)r&aQ(y)
 \end{pmatrix}.
\end{equation}
This is a linear map from $\Sym^4$ to itself for each fixed $y$.
On the space $\mathcal V=\Sym^4\times\Sym^4\times\R^3$, define
\begin{equation}\label{eq:p}
 p(X,Z,y)=\det\!\left((\det X)Z-\Phi_y(\adj X)\right),
 \qquad e=(I_4,I_4,0),
\end{equation}
where $\adj X$ denotes the classical adjugate. Each matrix entry inside
this determinant is homogeneous of degree five, and $p(e)=1$; hence
$p$ is homogeneous of degree twenty. We write $K=K(p,e)$.
The slice representation in Section~\ref{sec:slices} will in particular
verify that $p$ is hyperbolic with respect to $e$.

Identify $\mathcal V$ with $\R^{23}$ by listing the upper triangular
entries of $X$, then those of $Z$, each in lexicographic order of
$(i,j)$, followed by $y_1,y_2,y_3$.

\begin{theorem}\label{thm:main}
There exist fixed matrices $A_1,\ldots,A_{23},B_1,\ldots,B_{307}
\in\Sym^{100}$ such that, for every $x=(X,Z,y)\in\mathcal V$,
\[
 x\in K\quad\Longleftrightarrow\quad
 \text{there exists }v\in\R^{307}\text{ with }
 \sum_{\nu=1}^{23}x_\nu A_\nu+
 \sum_{j=1}^{307}v_jB_j\succeq0.
\]
The equivalence holds on the entire closed cone, including points
with singular $X$.
\end{theorem}

The sizes in Theorem~\ref{thm:main} are explicit bounds, with no
minimality assertion. The matrices are specified by a finite
coefficient rule in Section~\ref{sec:lift}.

The cone in \eqref{eq:p} was constructed in the companion
manuscript~\cite{parent}, whose Theorem~1.1 proves that it has no
homogeneous semidefinite representation in its original twenty-three
coordinates. This also excludes affine representations. Indeed, if a
cone containing $0$ equals $\{x:A_0+L(x)\succeq0\}$, where $L$ is
linear, then $A_0\succeq0$. For each point $x$ of the cone and $t>0$,
we have $A_0/t+L(x)\succeq0$; letting $t\to\infty$ gives $L(x)\succeq0$.
Conversely, $L(x)\succeq0$ implies $A_0+L(x)\succeq0$; thus the same
cone has the homogeneous representation $L(x)\succeq0$.
We use the companion result only for this comparison: the proof of
Theorem~\ref{thm:main} below is self-contained. The companion also
establishes hyperbolicity and, in Proposition~3.2, describes the part of
$K$ with $X\succ0$. Our slice calculation recovers those two facts as
part of the lift construction.

\subsection{Proof strategy}
The starting point is a family of planes in $\R^3$:
\[
 y=Y_\theta(\eta,\xi)
   =\eta(\cos\theta,\sin\theta,0)+\xi(0,0,1).
\]
They cover all values of $y$. On each plane, the matrix $Q(y)$ has a
Gram factorization linear in $(\eta,\xi)$ and trigonometric of degree
three in $\theta$. This gives a $20\times20$ linear pencil $L_\theta$
whose positive semidefinite locus is exactly the corresponding slice
of $K$, including its boundary.

The angular parameter does not itself furnish a finite linear matrix
inequality. We instead introduce a real linear functional on a
finite-dimensional space of trigonometric polynomials with coefficients
linear in the slice variables. We require it to be nonnegative on
$q^TL_\theta q$ for every vector $q$ of trigonometric polynomials of
degree at most two. These requirements form one finite matrix
inequality. Evaluation at a single angle provides a feasible functional
for every point of $K$.

This use of a finite matrix to encode positivity against a space of
test polynomials belongs to the moment and sum-of-squares approach to
semidefinite representations~\cite{lasserre2001,lasserre2009}.
Here exactness must hold at one fixed degree.

The main issue is the converse: an arbitrary feasible functional need
not be evaluation at an angle, or integration against a measure.
The degree-two matching identity in Section~\ref{sec:matching} supplies
enough test vectors to recover the matrix inequalities defining each
slice. Its construction rotates a fixed complex Gram factor by a
$2\times2$ unitary matrix. The rotation preserves the Gram matrix and
reduces a cubic trigonometric pairing to a first harmonic. Matching
its value and first derivative then identifies the pairing exactly.
This bounded-degree construction is the principal additional algebraic
ingredient of the lift. Section~\ref{sec:lift} applies it to prove
exactness for every feasible functional and completes the proof of
Theorem~\ref{thm:main}.
